% !TEX root = relatorio.tex
%
% Capítulo 9 — Atribuição, Remediação e Tutoria com IA
%
\chapter{Atribuição, Remediação e Tutoria com IA} \label{cap:desafios}

Este capítulo dá continuidade à descrição da geração de desafios apresentada no Capítulo~\ref{cap:geracao-desafios}, descrevendo o percurso do Learner desde a atribuição diária de um desafio até à sua resolução, passando pela remediação adaptativa e pelo modo de tutoria conversacional com IA quando a remediação não é suficiente. Na nomenclatura interna do sistema, os desafios são denominados \textit{tasks} e cada conceito de negócio tem correspondência direta a entidades de domínio bem delimitadas. A segurança aplicacional de ponta a ponta é tratada em separado no Capítulo~\ref{cap:seguranca}.

% -----------------------------------------------------------
\section{Atribuição aos Utilizadores} \label{sec:atribuicao}
% -----------------------------------------------------------

\subsection*{Entidades de Atribuição}

Dois agregados gerem o estado da aprendizagem por utilizador:

\begin{itemize}
    \item \textbf{\texttt{Enrollment}:} inscrição ativa de um utilizador num módulo, com o dia corrente do percurso, o estado (\texttt{ACTIVE}\,|\,\texttt{COMPLETED}\,|\,\texttt{EXPIRED}\,|\,\texttt{PAUSED}) e os pontos e \textit{streak} acumulados.
    \item \textbf{\texttt{TaskAssignment}:} atribuição concreta de um desafio a um utilizador num dia, com a permutação de opções aplicada, o prazo de 24~horas e o estado de submissão.
\end{itemize}
Esquema de campos completo no Apêndice~\ref{ap:detalhe-config} (Tabela~\ref{tab:model-assignment}).

\subsection*{Fluxo Pull-based}

A atribuição de desafios é \textit{pull-based}: o utilizador solicita a tarefa do dia via \texttt{GET /tasks/today}, e o serviço resolve quatro casos mutuamente exclusivos, por ordem de precedência: (1)~sessão de remediação adaptativa aberta; (2)~sessão de explicação ativa; (3)~assignment já existente para o \texttt{dayIndex} atual; (4)~criação de novo assignment.

No caso~(4), o sistema seleciona de forma determinística (por semente) uma das \texttt{TaskInstance} disponíveis (variedade de distratores) e a permutação das quatro opções (\texttt{optionOrder}) --- reprodutível para o mesmo utilizador, embora sem padrão previsível para quem não conhece a semente ---, e define \texttt{dueAt = now() + 24h}.

\subsection*{Permutação das Opções e Penalizações por Atraso}

A \texttt{optionOrder} é uma lista de índices (ex.: \texttt{[2,0,3,1]}) que define a ordem de apresentação no ecrã.
A resposta correta, sempre no índice~0 da instância, aparece numa posição aleatória para o utilizador.
A resposta submetida é um índice de posição de ecrã; o sistema resolve o índice canónico antes de validar, via \texttt{optionOrder[selectedOption]}.

A submissão após \texttt{dueAt} é aceite mas penalizada em função do atraso \textit{face ao próprio \texttt{dueAt}} (e não do tempo total desde a atribuição): 100\% dos pontos sem atraso, 75\% até 24\,h, 50\% até 72\,h, e 25\% (mínimo 1 ponto) além disso.

% -----------------------------------------------------------
\section{Remediação Adaptativa} \label{sec:consolidation}
% -----------------------------------------------------------

Quando um utilizador responde incorrectamente, o sistema não termina a interação: desencadeia automaticamente uma sessão de remediação adaptativa.
O objetivo é diagnosticar a causa do erro e apresentar três sub-tarefas \textit{scaffolded} que devolvam o utilizador ao nível da tarefa original.

\subsection*{Classificação do Padrão de Erro}

Antes de invocar a IA, o \texttt{ErrorPatternClassifier} classifica a falha de forma determinística em três categorias, com base no contexto da submissão:

\begin{table}[h!]
\centering
\caption{Padrões de erro classificados pelo ErrorPatternClassifier}
\label{tab:error-patterns}
\resizebox{\textwidth}{!}{%
\begin{tabular}{lll}
\hline
\textbf{Padrão} & \textbf{Condição} & \textbf{Estratégia de remediação} \\
\hline
\texttt{TIME\_PRESSURE} & Submissão após \texttt{dueAt} & Tarefas paceadas, passo a passo \\
\texttt{READING\_ERROR} & Opção seleccionada adjacente à correta (\texttt{|pos\_sel - pos\_corr| == 1}) & Instruções mais explícitas e claras \\
\texttt{WRONG\_CONCEPT} & Qualquer outro caso & Re-explicação do conceito de raiz \\
\hline
\end{tabular}%
}
\end{table}

\subsection*{Geração Assíncrona da Sessão}

Após classificar o erro, o \texttt{TaskService} persiste o resultado e delega ao \texttt{HandleStruggleService} via \texttt{triggerAsync}: a resposta HTTP é devolvida imediatamente ao utilizador, sem esperar pela geração do ramo adaptativo pelo LLM.
Em caso de falha ou \textit{timeout} do LLM, a sessão é marcada como \texttt{ABANDONED} e o assignment original é reposto em \texttt{PENDING}, e o utilizador nunca fica bloqueado.

\subsection*{Submissão e Resolução das Tarefas Adaptativas}

O utilizador acede às três sub-tarefas via \texttt{GET /struggle/enrollment/\{enrollmentId\}}.
Cada submissão pode resultar em quatro desfechos, avaliados após a conclusão de todas as tarefas: (1)~\textbf{ainda não concluídas}: devolver resultado parcial; (2)~\textbf{todas corretas}: marcar sessão como \texttt{RESOLVED} e repor assignment em \texttt{PENDING}; (3)~\textbf{falha, depth $<$ \texttt{MAX\_STRUGGLE\_DEPTH}}: gerar nova sessão de \textit{struggle} encadeada; (4)~\textbf{falha, depth $\geq$ \texttt{MAX\_STRUGGLE\_DEPTH} (= 1)}: criar sessão de explicação por IA.

A profundidade máxima de remediação está fixada em \texttt{MAX\_STRUGGLE\_DEPTH} = 1, uma decisão de desenho deliberada, assente em dois argumentos. O primeiro é pedagógico: as sub-tarefas partilham o formato de escolha múltipla do desafio original, pelo que encadear níveis adicionais do mesmo formato tem retorno decrescente: uma incompreensão que sobrevive a três sub-tarefas \textit{scaffolded} dificilmente se resolve com mais três no mesmo molde, e a tutoria conversacional (Secção~\ref{sec:explicacao-ia}) é o mecanismo desenhado precisamente para esse caso. O segundo é operacional: cada nível adicional multiplica as chamadas à API de IA e prolonga o tempo até o Learner regressar ao fluxo principal. O compromisso é reconhecido: tópicos avançados não dispõem de patamares intermédios antes do escalamento, e tornar a profundidade configurável por tópico fica identificado como trabalho futuro (Secção~\ref{sec:trabalho-futuro}).

As tarefas adaptativas atribuem zero pontos e não afetam o \textit{streak}: são modo de aprendizagem puro.

% -----------------------------------------------------------
\section{Tutoria com IA} \label{sec:explicacao-ia}
% -----------------------------------------------------------

Quando o utilizador esgota toda a profundidade de remediação sem sucesso, o sistema ativa o modo de tutoria conversacional com IA.
Esta sessão sintetiza o historial completo de erros e oferece uma explicação holística em prosa, seguida de um diálogo de esclarecimento sem limite de turnos. O acesso a este diálogo é limitado por utilizador, protegendo o custo das chamadas ao modelo de IA (detalhe no Capítulo~\ref{cap:seguranca}, Secção~\ref{sec:ratelimiting}).

\subsection*{Estados da Sessão de Explicação}

A sessão de explicação percorre três estados:
\[
\texttt{GENERATING} \;\to\; \texttt{ACTIVE} \;\to\; \texttt{RESOLVED}
\]

A transição \texttt{GENERATING} $\to$ \texttt{ACTIVE} ocorre quando a geração pelo LLM termina com sucesso.
Se a geração falhar ou esgotar o \textit{timeout} (60\,s via \texttt{withTimeoutOrNull}), a sessão transita diretamente para \texttt{RESOLVED} e o assignment é reposto em \texttt{PENDING} (\textit{fail-safe}: o utilizador nunca fica bloqueado).
A transição para \texttt{RESOLVED} pelo utilizador, via \texttt{POST /explanation/\{sessionId\}/resolve}, também repõe o assignment em \texttt{PENDING}, devolvendo-o ao fluxo diário normal.

\subsection*{Geração Holística e Diálogo Conversacional}

A geração holística é assíncrona (\texttt{generationExecutor}), permitindo que o utilizador navegue para o ecrã de explicação enquanto a geração ocorre em \textit{background}, com \textit{polling} até o estado ser \texttt{ACTIVE}. O contexto injetado inclui o historial completo de sessões de dificuldade, assemblado por joins relacionais sem pesquisa vectorial (Secção~\ref{sec:recuperacao} do Capítulo~\ref{cap:rag}).

Para o diálogo conversacional, cada mensagem é processada com o historial completo, usando a mesma estrutura multi-mensagem tipada descrita na Secção~\ref{sec:protecao-ia} do Capítulo~\ref{cap:ia} para isolar o input do utilizador.
A Figura~\ref{fig:seq-explanation-conv} (Apêndice~\ref{ap:detalhe-config}) detalha o fluxo completo desde a criação assíncrona até à resolução.
